\paragraph{\trLefschetzPencil{}。}\label{sub:lefschetz-pencils}
设 $\bar{X}$ 是一个 $(n+1)$ 维的\trSmooth{}\trProjectiveVariety，装备一个\trAmple{}\trLineBundle{}$\scr{L}$。取两个\trLinear{}无关的\trSection{}$s_0,s_\infty\in H^0(\scr{L})$，并考虑有理函数 $\bar{p}=s_0/s_\infty$，它的\trFiber{}是对 $z\in\CP^1$ 的诸\trHypersurface{}$\{s_0(y)=z\,s_\infty(y)\}$。$\bar{p}$ 的\trIndeterminacyLocus{}是\trBaseLocus{}$B$，它由 $s_0$ 与 $s_\infty$ 共同取零所定义。我们说 $\bar{p}$ 是一个\textbf{\trLefschetzPencil}（Lefschetz pencil），如果 $B$ 是\trSmooth{}的，并且 $\bar{p}$ 只有\trNondegenerate{}的\trCriticalPoint，且每根\trFiber{}中至多一个。我们还将额外假设 $s_\infty^{-1}(0)$ 本身是\trSmooth{}的，并考虑 $\bar{p}$ 在其补集上的限制，记作
\begin{equation}\label{eq:lefschetz-pencil-map}
p\colon X=\bar{X}\setminus s_\infty^{-1}(0)\longrightarrow\CC.
\end{equation}
在 $\scr{L}$ 上选一个\trMetric，使之给出一个\trKaehler{}\trForm。如\cref{thm:plurisubharmonic-examples}（i）中那样，$s_\infty$ 的对数范数定义了一个\trFiniteType{}的\trPlurisubharmonic{}函数 $\phi$。由此产生的设定与 §15 中引进的\trExact\trLefschetzFibration{}概念密切相关，但不完全相同。有两种处理这一出入的办法：一是论证 \eqref{eq:lefschetz-pencil-map} 的几何足够\trWellBehaved，使得必要的论证可以没有重大困难{}地照搬过来；或者，修改当前的局面，使之完全纳入先前设立的形式框架之中。

在第一种进路中，要点是证明 \eqref{eq:lefschetz-pencil-map} 上由 $\omega_\phi$ 给出的\trConnection{}具有\trWellBehaved{}的\trParallelTransport{}映射（尽管诸\trFiber{}非\trCompact）。为了看到这一点，沿\trBaseLocus{}作\trBlowup，得到 $\bar{p}$ 的\emph{图像}
\begin{equation*}
\widehat{X}=\{(z,y)\in\CP^1\times\bar{X}: s_0(y)=z\,s_\infty(y)\}.
\end{equation*}
它带有\trProjection{}$\hat{p}\colon\widehat{X}\to\CP^1$ 与 $\widehat{X}\to\bar{X}$。我们把给定的\trKaehler{}\trForm{}沿第二个映射\trPullback，并把结果记作 $\omega_{\widehat{X}}$。它不再是\trKaehler{}的，因为它沿\trExceptionalDivisor{}$\widehat{B}=\CP^1\times B$ \trDegenerate，但它确实限制为诸\trSmooth\trFiber{}上的一个\trKaehler{}\trForm，从而定义了一个辛\trConnection。这个\trConnection{}事实上在 $\widehat{B}$ 上是平凡的，所以它在补集 $\widehat{X}\setminus\widehat{B}$ 上的\trParallelTransport{}在远离\trSingularPoint{}处是\trWellDefined{}的。另一方面，如果我们从 $\hat{p}\colon\widehat{X}\setminus\widehat{B}\to\CP^1$ 中去掉 $\infty$ 处的\trFiber，就取回了 $p\colon X\to\CC$，其\trConnection{}由 $\omega_\phi$ 给出。

这个论证还多给出了一点信息，当放在\trGrading{}的背景中考虑时它格外有意思。也就是说，假设
\begin{equation}\label{eq:canonical-bundle-pencil}
\mathcal{K}_{\bar{X}}\cong\scr{L}^{\otimes-d}
\end{equation}
对某个 $d\in\ZZ$ 成立。$s_\infty^{\otimes-d}$ 在这一\trIsomorphism{}下的\trPreimage{}是一个有理的\trHolomorphic{}$(n+1)$-\trForm，它沿 $s_\infty^{-1}(0)$ 可以有零点（$d<0$）或极点（$d>0$）。于是，它将限制为 $X$ 上的一个\trHolomorphic{}\trVolume{}形式 $\eta_X$。在它——或者更确切地说，与之相伴的相对\trQuadraticForm{}$\eta_X^2/dz^2$——在场时，一切\trParallelTransport{}映射都自然地带上\trGrading，一如\cref{sub:relative-volume-forms}中那样。考虑我们的路径是一根充分大的\trCircle{}的情形。该\trFibration{}\trSmooth{}地延伸到 $\infty$ 之上，但相对\trQuadraticForm{}却不然。于是，绕无穷远一周的\trMonodromy{}（计其\trGrading）\trIsotopy{}于\trSuspension{}$S^k$（意即装备着常值\trGrading{}$-k$ 的恒同映射），对某个 $k\in 2\ZZ$。显然，$k$ 只依赖于 $\eta_X/dz$ 沿无穷远处\trFiber{}的极点阶，它等于 $d-2$。为简单起见，改为取一个平凡\trFibration{}$\CC^*\times Y\to\CC^*$，配以相对\trVolume{}形式 $z^{d-2}\eta_Y$。如果 $\lambda(t)=e^{2\pi it}$ 是单位\trCircle，那么沿 $\lambda|[0;t]$ 的\trParallelTransport{}是装备着\trGrading{}$(2d-4)t$ 的 $\mathrm{Id}_Y$，于是走完\trCircle{}一整圈之后，我们得到 $k=4-2d$。

我们眼下要讨论另一种进路的一个可能实现，即把 \eqref{eq:lefschetz-pencil-map} 变成 §15 意义下的一个\trExact\trLefschetzFibration。这需要在 $\CC$ 中选取一个基点 $*$；出于记号上的方便，取 $*=0$，因此假定它是 $p$ 的一个\trRegularValue。取定 $\CC$ 上的标准辛\trForm{}$\omega_\CC$ 及其标准原函数 $\theta_\CC$。稍后我们将想把给定的 $\omega_X$ \trDeformation{}为 $\omega'_X=\omega_X+rp^*\omega_\CC$，类似地把 $\theta_X$ \trDeformation{}为 $\theta'_X=\theta_X+rp^*\theta_\CC$，其中 $r\ge 0$。尽管缺乏\trCompactness，\trMoser{}论证在这里仍然适用，因为满足 $\omega^r(\cdot,Y^r)=\partial\theta^r_X/\partial r=p^*\theta_\CC$ 的\trVectorField{}$Y^r$ 处处（正向）正比于底空间 $\CC$ 上指向内部的\trRadial{}\trVectorField{}$-z\partial_z$ 的\trHorizontal{}\trLift。其结果是一族\trExact{}辛\trEmbedding
\begin{equation*}
\rho^r\colon X\longrightarrow X,\qquad (\rho^r)^*\omega^r_X=\omega_X,\qquad (\rho^r)^*\theta^r_X=\theta_X+d(\text{某个函数}).
\end{equation*}
我们发现，记住这些\trEmbedding{}的一个特殊性质是有用的，它可直接由我们对诸 $Y^r$ 的描述得出：
\begin{lemma}\label{thm:embedding-uniform-boundedness}
取一个\trCompact{}子集 $K\subset X$。那么存在另一个\trCompact{}子集 $\tilde{K}\subset X$，使得对一切 $r\ge 0$ 都有 $\rho^r(K)\subset\tilde{K}$。
\end{lemma}
\noqed
下一步是对给定的辛\trConnection{}使用\trRadial{}方向的\trParallelTransport，为 \eqref{eq:lefschetz-pencil-map} 构造无穷远处的\trTrivialization。它由一个开子集 $U\subset X$（使得 $p|(X\setminus U)$ 是\trProper{}的）连同 $U$ 与 $\CC\times U_0$ 之间的一个\trFibrewise{}\trDiffeomorphism{}组成，其中 $U_0=U\cap X_0$。在这个\trTrivialization{}中
\begin{equation*}
\theta_X=\theta_{X_0}+\kappa+dR,\qquad \omega_X=\omega_{X_0}+d\kappa,
\end{equation*}
其中 $\theta_{X_0}$、$\omega_{X_0}$ 是 $\theta_X$、$\omega_X$ 在 $X_0$ 上的限制；$\kappa$ 是某个在切于诸\trFiber{}的向量上取零的 1-\trForm，而 $R$ 是一个函数（两者实际上在所有满足 $z=0$ 的点 $(z,x)\in\CC\times U_0$ 处都取零值）。现在我们想修改\trParallelTransport，使之在我们给定的\trTrivialization{}中在无穷远处变得平凡。为此，取 $X_0$ 上一个在 $U_0$ 之外等于 $1$ 的\trCompactlySupported{}函数 $l$，并考虑修改后的形式
\begin{equation}\label{eq:modified-forms-infinity}
\tilde{\theta}_X=\theta_{X_0}+l\kappa+d(lR)+rp^*\theta_\CC,\qquad \tilde{\omega}_X=d\tilde{\theta}_X
\end{equation}
（对 $r\gg 0$），其理解是：在\trTrivialization{}适用的区域之外，它们被延拓成等于 $\theta'_X$、$\omega'_X$。取一个大的\trDisk{}$S\subset\CC$，并给它配上限制 $\theta_S=\theta_\CC|S$、$\omega_S=\omega_\CC|S$。“大”首先意味着 $S$ 的内部含有 $p$ 的一切\trCriticalValue；以后还会施加更多条件。取 $X_0$ 的一个\trCompact{}子集 $E_0$，它是某个\trPlurisubharmonic{}函数的一个大\trSublevelSet，并考虑 $p^{-1}(S)$ 中在无穷远处形如 $S\times E_0$ 的那块 $E$。这里的理解是：$E$ 将包含 $p^{-1}(S)$ 中在无穷远\trTrivialization{}里看不见的全部部分。此外，我们想选 $E_0$ 使得 $l$ 的\trSupport{}落在其内部。令 $\omega_E=\tilde{\omega}_X|E$、$\theta_E=\tilde{\theta}_X|E$，并考虑限制 $\pi=p|E\colon E\to S$。最后，取一个与 $\omega_E$ \trCompatible、在诸\trCriticalPoint{}附近与给定的 $I_X$ 一致、使 $\pi$ 成为\trPseudoholomorphic、并且在无穷远附近的\trTrivialization{}中看起来像自然乘积结构 $i\times I_{X_0}$ 的\trAlmostComplex{}结构 $I_E$。我们宣称，所得结果是一个\trExact\trLefschetzFibration。由于 $\tilde{\omega}_X$ 限制到诸\trFiber{}上时等于 $\omega_X$，且 $r$ 很大，显然 $\tilde{\omega}_X$ 仍是辛的，而且负\trLiouville\trVectorField{}沿 $\partial^v E=\pi^{-1}(\partial S)$ 指向内部；我们在考虑 \eqref{eq:kappa-deformation-family} 时见过这个论证。另一方面，我们知道负\trLiouville{}场沿 $E_0$ 的边界指向内部，故由构造，同样的结论沿整个 $\partial^h E$ 成立。事实上，我们已经把事情安排成使该\trFibration{}在 $\partial^h E$ 附近是平凡的，其意义见\cref{thm:original-definition-remark}。$\partial^h E$ 的弱 $I_E$-\trConvex{}性同样可以从 $E_0$ 的相应性质得出，而对 $\partial^v E$，它可以通过把\trPseudoholomorphic{}曲线\trProjection{}到底空间来证得。最后，假设 $\bar{X}$ 带有一个\trIsomorphism{} \eqref{eq:canonical-bundle-pencil}；由构造，$I_E$ 作为\trAlmostComplex{}结构\trHomotopic{}于 $I_X|E$，于是可以让 $\eta_X|E$ 通过这一\trHomotopy，得到 $(E,I_E)$ 的一个复\trVolume{}形式 $\eta_E$。

在验证完所需的性质之后，剩下要反思的是 \eqref{eq:modified-forms-infinity} 对\trSymplecticStructure{}所做的改动，以及这会如何影响\trFukayaCategory。$X$ 上如\cref{thm:plurisubharmonic-examples}（i）中那样给出的自然\trPlurisubharmonic{}函数 $\phi$ 是\trFiniteType{}的。取\trSublevelSet{}$X^{\phi,c}$，其中 $c$ 大于一切\trCriticalValue。把它取作\cref{thm:embedding-uniform-boundedness}中的 $K$，并考虑相应的 $\tilde{K}$。在构造中我们可以安排得使 $\tilde{\omega}_X=\omega^r_X$、$\tilde{\theta}_X=\theta^r_X$ 在 $\tilde{K}$ 上成立，从而有 $X^{\phi,c}$ 到 $E$ 之内的一个\trExact{}辛\trEmbedding。依\cref{thm:contact-type-embedding-quasi-iso}，$\scr{F}(E)$ 因此包含一个\trQuasiequivalent{}于 $\scr{F}(X^{\phi,c})$ 的全 \trAinfty-子\trCategory；这对我们的目的——用\trLefschetzFibration{}来做关于 $X$ 的陈述——已完全够用（$\scr{F}(E)$ 的任何其他\trObject{}都可以视为多余的而径直忽略）。如果不考虑\trTotalSpace{}$E$ 而是看\trFibration{}本身的几何，事情还会更简单。\trFiber{}$E_0$ 按定义是 $X_0$ 上一个\trPlurisubharmonic{}函数的大\trSublevelSet。$\pi\colon E\to S$ 的\trParallelTransport{}映射当然不是 $p\colon X\to\CC$ 相应映射的限制，但在 $E\setminus\partial_h E$ 的一大块\trCompact{}子集上（即无穷远处的\trTrivialization{}未曾触及的部分），它们将与那些限制一致。\footnote{译注：原书此处印作 $\partial_h E$（$h$ 为下标）；疑应为 $\partial^h E$（上标形；同页上文数处均作上标）；此处照原书排印，存疑备案。——译者注}可以把事情安排得使 $E_0$ 中每个\trLagrange\trSubmanifold{}都能被负\trLiouville{}流推入那一部分。这特别使我们能够照搬先前关于整体\trMonodromy{}的结果：如果 $\lambda$ 是 $S\setminus\mathrm{Critv}(\pi)$ 中沿边界正向平行而行的一根\trLoop，$L$ 是\trFiber{}$E_{\lambda(0)}$ 中的一个\trExact{}\trLagrange\trSubmanifold，那么\trMonodromy{}的像 $h_\lambda(L)$ 将\trIsotopy{}于 $L$。在 \eqref{eq:canonical-bundle-pencil} 在场时，这还有一个改进的版本：任何\trExactLagrangianBrane{}$L^{\#}$ 在\trMonodromy{}$h_\lambda$ 的自然\trGrading{}版本之下的像满足
\begin{equation}\label{eq:monodromy-graded-shift}
h_\lambda(L^{\#})\simeq S^{4-2d}L^{\#}.
\end{equation}
\begin{proposition}\label{thm:vanishing-cycles-split-generate}
设 $\pi\colon E^{2n+2}\to S$ 是一个来自\trLefschetzPencil{}的\trExact\trLefschetzFibration。假设它带有由 \eqref{eq:canonical-bundle-pencil} 所得的相对复\trVolume{}形式。取定基点 $*\in S$，并设 $\boldsymbol{V}^{\#}=(V^{\#}_1,\ldots,V^{\#}_m)$ 是\trFiber{}$M=E_*$ 中\trVanishingCycle{}的一个\trBasis，它们经任意选取的\trGrading{}（并且当 $n=1$ 而 $\mathrm{char}(\KK)\neq 2$ 时，取非平凡 \trPinStructure）被做成\trLagrange\trBrane。令 $\scr{A}=\Pi(Tw\scr{F}(M))$ 为\trTwistedComplex{}的 \trAinfty-\trCategory{}之\trSplitClosure，于是按定义 $H^0(\scr{A})=D^{\pi}\scr{F}(M)$。那么，只要 $d\neq 2$，$\boldsymbol{V}^{\#}$ 中的\trObject{}就\trSplitGeneration{}$\scr{A}$。
\end{proposition}
\begin{proof}
把 \eqref{eq:monodromy-graded-shift}、\cref{thm:twist-equals-tau} 与\cref{thm:shifted-twist-split-generates}结合起来。
\end{proof}
\begin{remark}\label{thm:d-equals-two-counterexample}
为了看到条件 $d\neq 2$ 确实是必要的，取 $\bar{X}=\CP^1\times\CP^1$，配以 $\scr{L}=\scr{O}_{\bar{X}}(1,1)$。此时 $X_0=\CC^*$，而两个\trVanishingCycle{}都\trIsotopy{}于同一个\trCircle{}$L=S^1\subset\CC^*$，带非平凡 \trPinStructure。把它记作 $L^{\#}$，而把带平凡 \trPinStructure{}的对应物记作 $L^{\#}\otimes\beta$。如果选取\trCoefficients{}域 $\KK$ 使其特征 $\neq 2$，那么依 \eqref{eq:pss-twisted-cohomology}，$HF^{*}(L^{\#},L^{\#}\otimes\beta)=H^{*}(L;|\beta|_{\KK})=0$。这表明 $L^{\#}\otimes\beta$ 不在由 $L^{\#}$ \trSplitGeneration{}的子\trCategory{}之中。

\begingroup\tolerance=9999\emergencystretch=3em\hbadness=10000
尽管如此，对 $d=2$ 的一般理论仍可以试着再往前推一点。写 $L_k=\tau_{V_{k+1}}\cdots\tau_{V_m}(L)$，并给它配备诱导的\trBraneStructure。诸\trExactTriangle
\[
\begin{tikzcd}
L^{\#}_k \ar[rr] && L^{\#}_{k-1} \ar[dl,"{[1]}"'] \\
& HF^{*}(V^{\#}_k,L^{\#}_k)\otimes V^{\#}_k \ar[ul] &
\end{tikzcd}
\]
中所有 $\rightarrow$ 映射的\trComposition{}是一个数 $c_{L^{\#}}\in HF^0(L^{\#}_m,L^{\#}_0)\cong HF^0(L^{\#},L^{\#})\cong H^0(L^{\#};\KK)=\KK$。细察\cref{thm:shifted-twist-split-generates}的证明即可看出，$c_{L^{\#}}=0$ 是使 $L^{\#}$ 含于由诸\trVanishingCycle{}\trSplitGeneration{}的子\trCategory{}的一个充分条件。另一方面，回想这些\trExactTriangle{}在\cref{thm:hf-tau-exact-triangle}中是如何构造的，就可以发现每个 $\rightarrow$ 映射都是一个\trInvariant，它对给定\trLefschetzFibration{}的一小块（只含一个\trCriticalPoint）的\trPseudoholomorphic\trSection{}进行计数。利用标准的\trGluing{}定理，可以把 $c_{L^{\#}}$ 写成一个与整个\trFibration{}相伴的\trInvariant{}$\Phi_{E/S,*}$，其针对的是实现\trIsotopy{} \eqref{eq:monodromy-graded-shift} 的移动\trLagrange\trBoundaryCondition{}$F$。\trEquivalence{}地，按原始的代数几何数据来表述：可以把 $F$ 想成 $X$ 的一个\trLagrange\trSubmanifold，它在 $\pi\colon X\to\CC$ 的底空间中一根大\trCircle{}上\trFibration，从而贴近无穷远处的\trFiber。$c_{L^{\#}}$ 随后对 $X$ 中边界含于 $F$ 的\trHolomorphicDisc{}进行计数。人们可以试着经一次到\trNormalCone{}的\trDegenerate（degeneration to the normal cone）或一次\trSymplecticCut（symplectic cut）来进一步分析这一局面。这使 $X$ \trDegenerate{}为两块的并——$\bar{X}$ 与\trTotalSpace{}$\scr{L}^{-1}\to s_\infty^{-1}(0)$——二者沿 $s_\infty^{-1}(0)$ \trGluing{}在一起。\trLagrange\trSubmanifold{}$F$ 移入第二块，于是可以把与 [79] 类似的技术应用于\trHolomorphicDisc{}的计数。我们在此不试图把这一点说得更显。\par\endgroup
\end{remark}
